\label{sec:related}

\subsection{Windows Event Log and Sysmon-Based Malware Detection}

Windows Event Logs have emerged as a primary data source for host-based malware detection
in IT environments. Achmad~\textit{et al.}\ \cite{achmad2025} leveraged Sysmon-enriched
Windows logs—capturing process creation, registry changes, and network connections—with
both supervised classifiers (Random Forest: F1~=~0.88) and unsupervised anomaly detection
(Local Outlier Factor: F1~=~0.99), demonstrating the complementarity of the two paradigms.
\"{O}nal and G\"{u}ven \cite{onal2025} conducted a systematic study of novel Windows
event types and their discriminative power for dynamic malware analysis, achieving
ROC-AUC of $0.962 \pm 0.009$; their feature taxonomy informs our event-ID selection.
Maniriho~\textit{et al.}\ \cite{maniriho2024} provided a comprehensive literature review
of Windows malware detection, identifying feature representation and class imbalance as
the two dominant open challenges—both addressed by this work. At the log-parsing level,
Drain \cite{he2017drain} enables scalable template extraction; DeepLog \cite{du2017deeplog}
frames anomaly detection as LSTM next-key prediction; LogAnomaly \cite{zhao2019loganomaly}
unifies sequential and quantitative anomaly signals. Our session-level feature engineering
captures the same sequential structure via Shannon entropy and time-delta statistics
without requiring deep model training, yielding substantially lower inference latency.
Brown~\textit{et al.}\ \cite{brown2018rnn} added attention mechanisms to recurrent networks
for interpretable anomaly localisation.

\subsection{Graph-Based and Provenance-Aware Detection}

A parallel line of research treats system events as vertices in a provenance graph rather
than independent feature vectors. Zhen~\textit{et al.}\ \cite{zhen2025} construct
heterogeneous graphs from audit logs and apply Graph Neural Networks (GNNs) to learn
structural dependencies among processes, files, and network endpoints. Han~\textit{et al.}\
\cite{han2020} proposed Unicorn, using online sketch-based provenance-graph summarisation
for runtime detection of Advanced Persistent Threats with low memory overhead. The spatial
awareness of these methods enables identification of complex attack chains invisible to
bag-of-features approaches. However, graph construction and traversal introduce significant
computational overhead; continuous graph updating in high-event-rate OT environments creates
latency incompatible with real-time detection requirements, motivating our choice of
lightweight session-level aggregation.

\subsection{Hybrid Static-Dynamic Analysis}

Ayub~\textit{et al.}\ \cite{ayub2023} proposed RWArmor, feeding static ML model outputs
into a dynamic behavioural analysis pipeline for early ransomware detection; accuracy
reaches 86--97\% within 30--120 second observation windows. Anderson and Roth
\cite{anderson2018} released EMBER, an open dataset for training static PE malware
classifiers, enabling reproducible feature-engineering studies. These hybrid approaches
improve detection fidelity for time-critical threats but depend on sandbox infrastructure
unavailable in operational OT environments where systems cannot be interrupted or diverted.

\subsection{AI-Driven Threat Hunting and Behavioural Analytics}

Beyond reactive classification, Chen~\textit{et al.}\ \cite{chen2022} combined ML, NLP,
and graph analytics for proactive threat hunting over large-scale event streams, framing
detection as discovery of anomalous behavioural narratives. Their system addressed the
``weak signal'' problem—where no single event is sufficient for classification, but a
sequence of events constitutes a detectable story. This narrative framing directly motivates
our session abstraction, which aggregates individual event signals into discriminative
session-level features. Their reliance on APT emulation frameworks for training data and the
high computational cost of combined NLP+graph processing remain open deployment challenges.

\subsection{ICS/SCADA Intrusion Detection}

Nafees~\textit{et al.}\ \cite{nafees2023} surveyed cyber-physical situational awareness
in smart grids, establishing the defining constraints of OT security: strict temporal
bounds ($<10$~ms decision latency in some scenarios), safety-critical operations that
cannot be interrupted, and tight coupling between cyber and physical processes. Carcano
\textit{et al.}\ \cite{carcano2011scada} proposed multidimensional critical-state
analysis for SCADA, encoding safe states as polytopes in the joint (measurement, command)
space. Goldenberg and Wool \cite{goldenberg2013modbus} built deterministic Modbus/TCP
models with zero false positives from protocol specifications. Goh~\textit{et al.}\
\cite{goh2016swat} released the SWaT dataset, still the most widely used ICS security
benchmark—capturing physical sensor and network measurements, not Windows host audit
trails. Linda~\textit{et al.}\ \cite{linda2009neural} demonstrated ML feasibility in
critical infrastructure. Mitchell and Chen \cite{mitchell2014survey} confirmed that
\emph{host-level} IDS methods are rare relative to network and process-level approaches.
Conti~\textit{et al.}\ \cite{conti2021survey} catalogued ICS testbeds and found no
publicly available Windows Event Log dataset for smart grid environments—the primary
data gap this work addresses.

\subsection{Machine Learning for Intrusion Detection}

Breiman's Random Forest \cite{breiman2001} consistently outperforms single classifiers on
tabular security features \cite{liao2013ids}. XGBoost \cite{chen2016xgboost} extends
gradient-boosted trees with column subsampling, achieving state-of-the-art results on
the KDD~Cup~1999 benchmark after duplicate removal \cite{tavallaee2009kdd}. Hink
\textit{et al.}\ \cite{hink2014machine} applied Random Forest to smart grid disturbance
and attack discrimination. Garcia-Teodoro~\textit{et al.}\ \cite{garcia2009anomaly}
surveyed anomaly-based IDS, establishing the statistical, knowledge-based, and ML
taxonomy this paper follows. Sommer and Paxson \cite{sommer2010outside} argued that
open-world operational environments undermine closed-world ML assumptions—their critique
motivates our multi-seed stability evaluation and concept drift discussion in
Section~\ref{sec:disc}. NIST SP~800-82 \cite{stouffer2015nist} defines the ICS security
control baseline against which we position detection and latency requirements.

\subsection{Research Gap and Positioning}

Table~\ref{tab:related} summarises the representative works reviewed above. A consistent
pattern emerges. Feature-based log methods \cite{achmad2025,onal2025} offer operational
efficiency but treat events as independent signals, missing multi-event attack narratives.
Graph-based models \cite{zhen2025,han2020} capture relational structure but incur latency
and preprocessing costs incompatible with OT real-time constraints. Hybrid pipelines
\cite{ayub2023,anderson2018} improve detection fidelity but require sandboxes. Threat
hunting systems \cite{chen2022} address context but impose high computational and data
requirements. All existing solutions are fundamentally \emph{IT-centric}: they assume
latency is negotiable, resources are abundant, and offline or sandbox-assisted analysis
is acceptable \cite{nafees2023}.

Smart grid control systems operate under a different reality. Decisions must be made
instantly; systems cannot be paused, replicated, or sandboxed; detection must function
within tight computational and temporal budgets. The result is a mismatch between
methodological sophistication and operational feasibility that no prior work resolves.
This paper addresses the gap by proposing a lightweight, session-level feature engineering
pipeline on Windows Event Logs, explicitly designed for latency-sensitive OT environments.

\begin{table*}[!t]
  \caption{Comparison of Representative Related Works}
  \label{tab:related}
  \centering
  \footnotesize
  \begin{tabular}{llllll}
    \hline
    Reference & Method & Strength & Key Limitation & Data Source & Env. \\
    \hline
    Achmad et al.\ \cite{achmad2025}        & Sysmon + ML             & Practical, low overhead         & No OT/ICS context     & Sysmon logs        & IT \\
    \"{O}nal \& G\"{u}ven \cite{onal2025}   & Windows events + ML     & Rich feature taxonomy           & Sandbox dependency    & Windows logs       & IT \\
    Zhen et al.\ \cite{zhen2025}            & GNN on audit logs       & Relational attack modeling      & High latency          & Audit logs         & IT \\
    Han et al.\ \cite{han2020}              & Provenance graph (APT)  & Online, memory-efficient        & Scalability limits    & System calls       & IT \\
    Ayub et al.\ \cite{ayub2023}            & RWArmor hybrid          & Early ransomware detection      & Sandbox bottleneck    & Behavioral logs    & IT \\
    Anderson \& Roth \cite{anderson2018}    & Static PE features      & Open reproducible dataset       & Static only           & PE binaries        & IT \\
    Chen et al.\ \cite{chen2022}            & ML+NLP+Graph hunting    & Narrative-level context         & High compute cost     & Event streams      & IT \\
    Nafees et al.\ \cite{nafees2023}        & Smart grid survey       & Domain-specific OT constraints  & No implementation     & OT systems         & OT \\
    \hline
    \textbf{This work}                      & RF on WEL sessions      & Lightweight, OT-deployable      & Synthetic data only   & Windows Event Logs & OT \\
    \hline
  \end{tabular}
\end{table*}
